% Gerado por ferramentas/importar_relatorios.py a partir de CWRU/RELATORIO_CWRU-10.11.2026/RELATORIO_CWRU.tex.
% Revisado à mão na dissertação (versão enxuta): edite este arquivo, não a fonte.
\subsection{Base de dados}
\label{cwru:sec:dados}

Foram usados os sinais do acelerômetro do lado do acionamento (DE), a 48\,kHz, dos conjuntos \textit{48k Drive End Bearing Fault Data} e \textit{Normal Baseline Data} \cite{cwru}. O rolamento é um 6205-2RS JEM SKF, com frequências de defeito de 5,4152 (pista interna, BPFI), 3,5848 (pista externa, BPFO) e 2,3570 (esfera, BSF) vezes a frequência de rotação. O subconjunto tem delineamento cruzado completo: máquina sã (N) e defeitos por eletroerosão na pista interna (IR), na esfera (B) e na pista externa na zona de carga (OR@6h), com diâmetros de 0,007, 0,014 e 0,021~polegada, nas cargas de 0 a 3\,HP. São 40 gravações de cerca de 10\,s. As inconsistências conhecidas dos arquivos foram tratadas explicitamente: arquivos com variáveis de outras gravações (usou-se sempre a de número igual ao do arquivo), um arquivo com apenas 1,3\,s, uma rotação ausente (usou-se a nominal) e cinco gravações com saturação em \(\pm 3{,}418\,g\), em fração desprezível. Cada gravação foi dividida em janelas de 0,25\,s (cerca de 7 rotações), sem sobreposição, num total de 1\,376 janelas: 140 da máquina sã, 396 de IR, 420 de B e 420 de OR.

\subsection{Métodos}
\label{cwru:sec:metodos}

O extrator reproduz a lógica do MOMEA sobre \(|x - \bar{x}|\) (início no primeiro cruzamento, término após o DDT, HLT re-armável com ROVL e truncamento por MAXDUR), com os tempos escalados para 48\,kHz: DDT e HLT de 1\,ms e duração máxima de 20\,ms. Foram comparados dois limiares: \textbf{fixo}, de 0,3\,\(g\), como o MOMEA opera hoje; e \textbf{flutuante}, de 1,5 vez o RMS da janela anterior (causal), com mínimo de 0,3\,\(g\), análogo ao de sistemas comerciais de EA \cite{vallen2022} e implementável na lógica. O fator 1,5 e o DDT foram escolhidos por varredura, pelo critério de a ordem do intervalo entre \textit{hits} coincidir com BPFO e BPFI; o critério usa os rótulos, mas não a divisão de treino e teste de nenhuma validação.

Três grupos de descritores foram calculados por janela:
\begin{itemize}
	\item \textbf{clássicos} (8): RMS, pico, fator de crista, curtose, assimetria e fatores de forma, de impulso e de folga;
	\item \textbf{envelope} (7): curtose do envelope e amplitudes do espectro do envelope (filtro de 1 a 10\,kHz) em 1\(\times\) e 2\(\times\) BPFO, BPFI e 2\(\times\)BSF, em dB acima da mediana \cite{randall2011};
	\item \textbf{\textit{hits}} (14): taxa, frações truncadas e com ROVL, percentis dos descritores e, como medidas de periodicidade, a ordem do intervalo entre \textit{hits}, \((1/\mathrm{mediana}\,\Delta t)/f_r\), e seu coeficiente de variação.
\end{itemize}

A MLP tem os hiperparâmetros do estudo da bomba \cite{pedregosa2011}, e uma \textit{Random Forest} serve de referência. As tarefas são quatro classes (N, IR, B, OR) e dez classes (tipo e diâmetro), com três validações de independência crescente: \textbf{janelas sorteadas} (5 partes estratificadas); \textbf{carga deixada de fora} (LOLO, 4 partes), que separa gravações, mas não rolamentos; e \textbf{diâmetro deixado de fora} (LOSO, 3 partes, só em quatro classes), em que o teste contém rolamentos físicos nunca vistos, análogo ao LORO da bomba. As métricas são a acurácia por janela e por gravação (voto majoritário).
